\section{Planteamiento del problema}
    Cada año cerca de un millón de estudiantes presentan la prueba Saber 11 en
    el calendario A. El puntaje global que obtienen, en una escala de 0 a 500,
    resume su desempeño en cinco pruebas, lectura crítica, matemáticas,
    competencias sociales y ciudadanas, ciencias naturales e inglés y es el 
    criterio que utilizan las universidades, el programa \textit{Generación E} 
    y otrosmecanismos de acceso a la educación superior \cite{icfes2022guia}.

    \vspace{10pt}
    
    Dado que el ICFES recoge al momento de la inscripción un cuestionario sobre 
    el hogar del estudiante y conoce las características del colegio en el que 
    cursó la media académica, es posible preguntarse en qué medida ese puntaje 
    es predecible a partir de información que existe \textit{antes} de presentar 
    el examen y qué factores pesan más en esa predicción.

    \subsection{Pregunta de investigación}
        ¿En qué medida las características socioeconómicas del hogar, del colegio 
        y aspectos demográficos del estudiante permiten predecir el puntaje global 
        de la prueba Saber 11 y cuáles de esas características tienen mayor peso?

    \subsection{Objetivos}
        \subsubsection{Objetivo general}
            Construir y evaluar un modelo de regresión que prediga el puntaje
            global de la prueba Saber 11 del periodo 2022-4 a partir de variables
            socioeconómicas, familiares y del colegio e interpretar la
            contribución de cada grupo de variables.

            \vspace{10pt}

        \subsubsection{Objetivos específicos}
            \begin{enumerate}
                \item Describir el conjunto de datos \textit{Resultados únicos
                Saber 11} (origen, licencia, variables, tamaño y calidad) y
                caracterizar la distribución de la variable objetivo.

                \item Identificar, mediante análisis exploratorio, las variables
                con mayor asociación con el puntaje global y las que deben
                excluirse por fuga de información o por falta de valor predictivo.

                \item Preparar los datos (tratamiento de valores faltantes,
                codificación de variables ordinales y nominales, ingeniería de
                variables) y establecer un modelo de referencia con su métrica
                inicial.

                \item Entrenar y comparar al menos dos modelos adicionales al de
                referencia, ajustar sus hiperparámetros con validación cruzada y
                evaluarlos con métricas apropiadas para regresión (MAE, RMSE y
                $R^2$).

                \item Interpretar el modelo final (importancia de variables y
                análisis de errores) y comunicar los resultados mediante
                visualizaciones, un informe técnico y una sustentación.
            \end{enumerate}

            Este informe cubre los objetivos específicos 1 y 2, que corresponden
            a la Fase 1 del proyecto, los restantes se abordan en las fases 2, 3 y
            4.

    \subsection{Hipótesis de trabajo}
        \begin{itemize}
            \item \textbf{H1:} El nivel socioeconómico del hogar, medido por el
            estrato, la educación de los padres y el acceso a internet y
            computador es el grupo de variables con mayor asociación con el
            puntaje global.

            \item \textbf{H2:} Las características del colegio, su naturaleza
            oficial o privada, su área urbana o rural, su jornada y su calendario
            aportan información adicional a la del hogar.

            \item \textbf{H3:} Existe una brecha regional, de modo que el
            departamento donde se ubica el colegio explica parte de la variación
            del puntaje incluso después de considerar el hogar y el colegio.
        \end{itemize}

    \subsection{Alcance y restricciones}
        El problema se formula como una regresión sobre el puntaje global. Los
        puntajes por área y el nivel de desempeño en inglés no pueden usarse como
        variables explicativas, porque el puntaje global se calcula a partir de
        ellos y su inclusión constituiría una fuga de información. Se utiliza un 
        único periodo, el 2022-4, para no mezclar escalas ni cuestionarios de 
        distintos años y para mantener el volumen manejable en el entorno gratuito 
        de Google Colab, adicionalmente, se realiza elanálisis es correlacional
        que identifica asociaciones útiles para predecir y no relaciones causales.